\documentclass[12pt,a4paper]{article}
\usepackage{fontawesome5}
\providecommand{\faPencilAlt}{\faIcon{pencil-alt}}
\usepackage{styles/couleurs}
\usepackage{styles/mise_en_page}
\usepackage{styles/boites}
\usepackage{styles/maths}
\usepackage{styles/figures}
\usepackage{enumitem}
\usepackage{array}
\usepackage{tabularx}
\usepackage{amsmath}
\usepackage{amssymb}
\usepackage{tikz}
\usepackage{xparse}
\usetikzlibrary{calc,patterns,arrows.meta}
\setlist[itemize]{leftmargin=1.4em,itemsep=0.15em,topsep=0.2em}
\setlist[enumerate]{leftmargin=*,itemsep=0.25em,topsep=0.25em}
\newcommand{\lignepointillee}{\noindent\dotfill\par\vspace{0.45em}}
\newcommand{\lignesreponse}[1]{%
\begin{tcolorbox}[enhanced,colback=white,colframe=NeutreC!40,boxrule=0.5pt,
arc=1.5mm,left=3mm,right=3mm,top=2mm,bottom=1mm,before skip=2pt,after skip=2pt]
\foreach \i in {1,...,#1}{\lignepointillee}
\end{tcolorbox}}
\newcommand{\titreseq}[1]{%
\begin{center}
{\sffamily\LARGE\bfseries\color{TitreC}Séquence 4 -- Calcul littéral}\\[0.2em]
{\sffamily\large\bfseries\color{SavoirC}#1}
\end{center}}
\begin{document}

\titreseq{Activité 4 -- Simplifier l'écriture d'une expression}

\section*{Partie A -- Des écritures différentes}

Associe les écritures qui désignent le même nombre.

\[
3\times x \qquad x+x+x \qquad 3x
\]

\[
x\times x \qquad x^2
\]

\[
1\times x \qquad x
\]

\[
x\times5 \qquad 5x
\]

\lignesreponse{4}

\section*{Partie B -- Conventions}

\begin{tcolorbox}[colback=SavoirC!5,colframe=SavoirC!60,boxrule=0.7pt,arc=2mm]
On peut ne pas écrire le signe \(\times\) devant une lettre ou une parenthèse :

\[
3\times x=\dotfill
\qquad
5\times(x+2)=\dotfill
\]

On écrit le nombre avant la lettre :

\[
x\times7=\dotfill
\]

On écrit :

\[
x\times x=\dotfill
\qquad
1\times x=\dotfill
\]
\end{tcolorbox}

\section*{Partie C -- Simplifier}

\[
A=8\times x
\qquad
B=x\times4
\qquad
C=2\times x\times x
\]

\[
D=1\times x+5
\qquad
E=3\times(x+7)
\qquad
F=x\times x\times x
\]

\lignesreponse{5}

\section*{Partie D -- Ne pas confondre}

Explique pourquoi \(3+x\) et \(3x\) ne désignent pas la même chose.
Choisis une valeur de \(x\) pour le montrer.

\lignesreponse{4}

\section*{Partie E -- Remplacer une lettre}

Calcule la valeur de \(A=3x-5\) pour \(x=4\), puis pour \(x=-2\).
Écris les parenthèses nécessaires.

\lignesreponse{5}

\end{document}
